\section*{Chapter \ref{ch:b2:energy-geology} --- Energy Geology}
% numbers from the Book 4 second-writer scratch calculator calc/ch30.py

\begin{solution}{exo:b2:energy-geology:1}
(a) 1.58: type~I; 0.95: type~III; 1.24: type~II; 1.12: mixed type~II/III. The type~III kerogen, at
0.95, yields mainly gas. (b) H/C $= 150/100 = 1.50$, O/C $= 8/100 = 0.08$. (c) Per \qty{100}{g}:
$77.4/12.011 = \qty{6.444}{mol}$ C, $9.8/1.008 = \qty{9.722}{mol}$ H, $8.2/15.999 =
\qty{0.513}{mol}$ O; H/C $= 1.51$ and O/C $= 0.0795$, in agreement with the formula.
\end{solution}

\begin{solution}{exo:b2:energy-geology:2}
(a) 0.45: immature, below the onset at \qty{0.60}{\percent}; 0.68: generating oil; 1.1: still in the
oil window, near its end; 1.6: past 1.20--\qty{1.30}{\percent}, so oil breaks down to methane and the
rock yields gas; 2.5: gas only, generation continuing up to about \qty{4.8}{\percent}. (b) No: methane
generation from type~III begins near \qty{0.73}{\percent}. (c) $F = (\ln R_o + 1.6)/3.7$: $(-0.511 +
1.6)/3.7 = 0.294$ and $(0.262 + 1.6)/3.7 = 0.503$.
\end{solution}

\begin{solution}{exo:b2:energy-geology:3}
(a) In the reservoir, $6 \times 10^6 \times 25 \times 0.18 \times 0.70 = \qty{1.89e7}{m^3}$;
$\mathrm{STOIIP} = 1.89 \times 10^7/1.25 = \qty{1.51e7}{m^3}$. (b) $0.30 \times 1.512 \times 10^7 =
\qty{4.54e6}{m^3}$ and $0.40 \times 1.512 \times 10^7 = \qty{6.05e6}{m^3}$.
\end{solution}

\begin{solution}{exo:b2:energy-geology:4}
(a) $V = 10^7 \times 200 = \qty{2e9}{m^3}$ and $q_R = 2.5 \times 10^6 \times 2 \times 10^9 \times 60
= \qty{3.0e17}{J}$. (b) Below \qty{90}{\celsius}: a \omterm{def:b2:energy-geology:geothermal}{low-enthalpy resource}, for direct use of the
heat. (c) $P = 0.25 \times 3.0 \times 10^{17}/(30 \times \qty{3.156e7}{s}) = \qty{79.2}{MW}$ of heat.
\end{solution}

\begin{solution}{exo:b2:energy-geology:5}
(a) $k_2/k_1 = \exp[(E/R)(1/T_1 - 1/T_2)] = 2$ with $1/373.15 - 1/383.15 = \qty{6.994e-5}{K^{-1}}$
gives $E = 8.314 \times \ln 2/(6.994 \times 10^{-5}) = \qty{82.4}{kJ/mol}$. (b) $\exp(26800 \times
6.994 \times 10^{-5}) = 6.52$. (c) The doubling rule, whose effective energy is lower, gives
relatively more weight to the long times spent at low temperature; for a type~I kerogen with a high
activation energy, the TTI announces generation too early.
\end{solution}

\begin{solution}{exo:b2:energy-geology:6}
(a) $x = 1/2$ in the approximation gives $(ART^2/\beta E)\,e^{-E/RT} = \ln 2$; taking logarithms,
$E/RT = \ln[AT^2/(\beta(E/R)\ln 2)]$, and the stated form follows. Starting from any guess near
\qty{400}{K} and iterating with $\beta$ in kelvins per second: at \qty{1}{\celsius\per\Myr}
($\beta = \qty{3.169e-14}{K/s}$), $T = \qty{409.6}{K}$, \qty{136.4}{\celsius}; at \qty{10}{K} per
minute ($\beta = \qty{0.1667}{K/s}$, $\num{5.26e12}$ times faster), $T = \qty{719.0}{K}$,
\qty{445.9}{\celsius}. (b) The difference is \qty{309.5}{K}; the estimate $6.259 \times
\ln(5.26 \times 10^{12}) = 6.259 \times 29.29 = \qty{183}{K}$ is far too small, because $RT^2/E$
itself grows with $T$ (to \qty{19.3}{K} at \qty{719}{K}): the linear shift holds only for small
changes of rate.
\end{solution}

\begin{solution}{exo:b2:energy-geology:7}
$\lambda_5 - \lambda_8 = \qty{8.29725e-10}{yr^{-1}}$ and $r_0 = 0.7204/99.2742 = 0.0072567$.
(a) At \qty{2.0}{\Ga}, $r = 0.0072567\,e^{1.6595} = 0.03814$, share $0.03814/1.03814 =
\qty{3.67}{\percent}$; at \qty{1.8}{\Ga}, $r = 0.03231$, share \qty{3.13}{\percent}. (b) $r =
3/97 = 0.030928$ and $t = \ln(0.030928/0.0072567)/(8.29725 \times 10^{-10}) = \qty{1.75}{\Ga}$.
(c) Both the depleted ore and normal uranium have decayed since by the same factors, so the ratio of
their $^{235}$U shares is the same today as just after the reactions: the ore then held half the
normal share, so about half of its $^{235}$U had been fissioned.
\end{solution}

\begin{solution}{exo:b2:energy-geology:8}
(a) $10^9/1.872 = \qty{5.34e8}{m^3}$; $10^9/621.4 = \qty{1.61e6}{m^3}$; $10^9/686.4 =
\qty{1.46e6}{m^3}$. (b) $7 \times 10^9/583 = \qty{1.20e7}{m^3}$ of \ce{CO2}, filling
$0.35 \times 0.30 = 0.105$ of the rock volume: $1.20 \times 10^7/0.105 = \qty{1.14e8}{m^3}$ of rock,
a mean thickness of $1.143 \times 10^8/(5 \times 10^6) = \qty{22.9}{m}$.
\end{solution}

\begin{solution}{exo:b2:energy-geology:9}
(a) $V = \qty{9e9}{m^3}$ and $q_R = 2.6 \times 10^6 \times 9 \times 10^9 \times 235 =
\qty{5.50e18}{J}$. (b) Over $30 \times \qty{3.156e7}{s} = \qty{9.47e8}{s}$: $R_g = 0.1$ gives
\qty{581}{MW} and $0.25$ gives \qty{1450}{MW} of heat. (c) $6 \times (10 \text{ to } 20) = 60$ to
\qty{120}{MW} of electricity, $60/581 = 0.10$ to $120/581 = 0.21$ of the thermal power: most of the
heat is rejected, as in any power plant working with a modest temperature difference.
\end{solution}

\begin{solution}{exo:b2:energy-geology:10}
(a) The strength is $0.6 \times (80 - 39) = \qty{24.6}{MPa} > \qty{22}{MPa}$: stable. Failure needs
$0.6\,\Delta p = 24.6 - 22$, $\Delta p = \qty{4.33}{MPa} = \qty{43}{bar}$, of the same order as the
\qty{32}{bar} that bounded the Denver earthquakes. (b) $(24.6 - 24)/0.6 = \qty{1.0}{MPa} =
\qty{10}{bar}$: a fault already close to failure needs only a small rise. (c) Injected fluid raises
the pressure in a zone that spreads by diffusion through permeable rock and fault zones, kilometres
away over years; wherever it meets a fault already near failure, a modest rise is enough.
\end{solution}

\begin{solution}{exo:b2:energy-geology:11}
(a) Write $x = e^{-u}$ with $u = (ART^2/\beta E)\,e^{-E/RT}$. A transformation ratio of
\qty{10}{\percent} means $x = 0.9$, $u = \ln(10/9)$; \qty{90}{\percent} means $u = \ln 10$. As
$\mathrm{d}\ln u/\mathrm{d}T = E/RT^2 + 2/T \approx E/RT^2$, $\Delta(\ln u) \approx (E/RT^2)\,\Delta T$,
which gives the formula, with $\ln\ln 10 - \ln\ln(10/9) = 0.834 + 2.250 = 3.084$. (b) With $T =
\qty{409.6}{K}$, $RT^2/E = \qty{6.26}{K}$ and $\Delta T = \qty{19.3}{K}$; the exact single-energy
curve gives \qty{18.6}{K}, from 125.3 to \qty{143.8}{\celsius}. (c) The solid curve rises from
\qty{10}{\percent} near \qty{122}{\celsius} to \qty{90}{\percent} near \qty{148}{\celsius}, over about
\qty{26}{K}: the spread of activation energies widens the conversion interval.
\end{solution}

\begin{solution}{exo:b2:energy-geology:12}
(a) $2 \times 10^{10}/20 = \qty{1.0e9}{m^3}$ (\qty{1.0}{km^3}) and $2 \times 10^{10}/60 =
\qty{3.3e8}{m^3}$. (b) $2 \times 10^{10}/621.4 = \qty{3.22e7}{m^3}$. (c) $3.22 \times 10^7/0.02 =
\qty{1.61e9}{m^3}$ of pores, which with $50 \times 0.30 = \qty{15}{m}$ of pore space per square
metre requires $1.61 \times 10^9/15 = \qty{1.07e8}{m^2}$, about \qty{107}{km^2}.
\end{solution}

\begin{solution}{pb:b2:energy-geology:1}
\noindent\textbf{Part I.}
\textbf{1.} At \qty{70}{\Ma}: $50 \times 60 = \qty{3000}{m}$, $15 + 25 \times 3.0 =
\qty{90}{\celsius}$. Today: $3000 + 70 \times 30 = \qty{5100}{m}$, \qty{142.5}{\celsius}.
\textbf{2.} $60 \times 0.025 = \qty{1.5}{\celsius\per\Myr}$, then $30 \times 0.025 =
\qty{0.75}{\celsius\per\Myr}$.
\textbf{3.} 15--\qty{20}{\celsius}: \qty{3.33}{\Myr}; each interval from 20 to \qty{90}{\celsius}:
$10/1.5 = \qty{6.67}{\Myr}$; each from 90 to \qty{140}{\celsius}: $10/0.75 = \qty{13.3}{\Myr}$;
140--\qty{142.5}{\celsius}: \qty{3.33}{\Myr}.
\textbf{4.} Interval values $\Delta t_n 2^n$, with running sums: 15--20 ($n = -9$) 0.0065; 20--30
0.026 (0.033); 30--40 0.052 (0.085); 40--50 0.104 (0.189); 50--60 0.208 (0.397); 60--70 0.417
(0.814); 70--80 0.833 (1.65); 80--90 1.67 (3.31); 90--100 6.67 (9.98); 100--110 13.3 (23.3);
110--120 26.7 (50.0); 120--130 53.3 (103); 130--140 107 (210); 140--142.5 ($n = 4$) 53.3: TTI
today 263.
\textbf{5.} 15 is passed in the interval 100--\qty{110}{\celsius}, at $100 + 10 \times
(15 - 9.98)/13.33 = \qty{103.8}{\celsius}$, reached $(103.8 - 90)/0.75 = \qty{18.4}{\Myr}$ after
\qty{70}{\Ma}: at \qty{51.6}{\Ma} and $(103.8 - 15)/0.025 = \qty{3550}{m}$.
\textbf{6.} 75: \qty{124.7}{\celsius}, \qty{23.7}{\Ma}, \qty{4390}{m}. 160: \qty{135.3}{\celsius},
\qty{9.6}{\Ma}, \qty{4810}{m}.
\textbf{7.} No: at 263 it is past the end of oil generation (160), since about \qty{10}{\Ma}.

\noindent\textbf{Part II.}
\textbf{8.} $k = 9.21 \times 10^{13}\,e^{-26800/363.15} = \qty{8.20e-19}{s^{-1}}$. The rock was
below \qty{90}{\celsius} for the \qty{50}{\Myr} before \qty{70}{\Ma}, so the converted fraction is
less than $k\,t = 8.20 \times 10^{-19} \times 50 \times 3.156 \times 10^{13} = \num{1.29e-3}$.
\textbf{9.} $\beta = 0.75/(3.156 \times 10^{13}) = \qty{2.376e-14}{K/s}$, $T = \qty{415.65}{K}$:
$ART^2/(\beta E) = \num{2.498e28}$ and $e^{-E/RT} = \num{9.95e-29}$, so $u = 2.486$, $x = 0.0833$ and
the transformation ratio is 0.917.
\textbf{10.} \qty{134.7}{\celsius}.
\textbf{11.} $70 - (134.7 - 90)/0.75 = \qty{10.4}{\Ma}$.
\textbf{12.} The TTI places peak generation at \qty{23.7}{\Ma}, the kinetics \qty{13}{\Myr} later,
because the doubling rule corresponds to $E \approx \qty{82}{kJ/mol}$ and weights the low
temperatures too much for a kerogen of \qty{223}{kJ/mol} (exercise~\ref{exo:b2:energy-geology:5}).
\textbf{13.} Earlier: a lower energy lowers the temperature of half conversion.

\noindent\textbf{Part III.}
\textbf{14.} Rock: $4 \times 10^8 \times 50 \times 2400 = \qty{4.8e13}{kg}$; potential $0.030
\times 4.8 \times 10^{13} = \qty{1.44e12}{kg}$; generated $0.92 \times 1.44 \times 10^{12} =
\qty{1.32e12}{kg}$.
\textbf{15.} $1.3248 \times 10^{12}/850 = \qty{1.56e9}{m^3}$.
\textbf{16.} $1.2 \times 10^7 \times 40 \times 0.22 \times 0.75/1.30 = \qty{6.09e7}{m^3}$.
\textbf{17.} $6.09 \times 10^7/1.56 \times 10^9 = 0.039$, \qty{3.9}{\percent}: most generated oil
stays in the source rock, is lost along the migration paths or escapes to the surface.
\textbf{18.} $0.35 \times 6.09 \times 10^7 = \qty{2.13e7}{m^3}$.
\textbf{19.} About two thirds of the oil stays in the reservoir after primary and secondary
recovery; \ce{CO2} injection could recover a further 10 to \qty{15}{\percent} of the oil in place.

\noindent\textbf{Part IV.}
\textbf{20.} $15 + 25 \times 2 = \qty{65}{\celsius}$ and $0.101 + 1000 \times 9.807 \times 2000/10^6
= \qty{19.7}{MPa}$: both above the critical point (\qty{31.0}{\celsius}, \qty{7.38}{MPa}), so the
\ce{CO2} is supercritical, at about \qty{686}{kg.m^{-3}}.
\textbf{21.} $2.13 \times 10^7 \times 1.30 = \qty{2.77e7}{m^3}$.
\textbf{22.} Pore volume $1.2 \times 10^7 \times 40 \times 0.22 = \qty{1.056e8}{m^3}$, and $E =
2.77 \times 10^7/1.056 \times 10^8 = 0.262$, far above the 1 to \qty{4}{\percent} of saline
formations: the produced oil has emptied a known, sealed space.
\textbf{23.} $6 \times 10^8 \times 80 \times 0.20 \times 691 \times 0.02 = \qty{1.33e11}{kg}$,
\qty{133}{Mt}.
\textbf{24.} At \qty{1}{Mt} a year, about 133 years.
\textbf{25.} $M = \rho\,N_p\,B_o = 686 \times 2.77 \times 10^7 = \qty{1.90e10}{kg}$: the depleted
reservoir can store \textbf{\qty[text-series-to-math=true]{19.0}{Mt}} of \ce{CO2}.
\end{solution}
